\section{Structured Abstractions for Open-World Agents}
\label{sec:abstractions}

The preceding sections established the setting (open-world environments), the agent's internal structure (hybrid planning and learning), and the problem (novelty as mismatch with an associated taxonomy). This section addresses the question that connects them: given that novelty is unavoidable, what determines whether an agent can respond to it effectively?

The central argument of this dissertation is that the answer lies in \emph{structured abstractions}---representations that organize the agent's knowledge and behavior at appropriate levels of granularity, exposing the compositional structure of tasks and interactions. We develop this argument by first situating the use of abstraction in this dissertation within the broader literature and then examining two specific forms of abstraction in the subsequent chapters.

\subsection{Abstraction in Agent Systems}
\label{sec:abstraction_definition}

Abstraction is a foundational concept in artificial intelligence and reinforcement learning. At its core, an abstraction is a mapping that reduces the complexity of a representation while preserving structure relevant to a particular purpose. \citet{konidaris2019necessity} argues that the role of abstraction is to capture the complexity of the \emph{task} rather than the complexity of the \emph{agent}: a well-chosen abstraction strips away details of the agent's embodiment and sensory apparatus, exposing the underlying structure of the problem to be solved. This perspective is central to the approach taken in this dissertation.

In the reinforcement learning literature, two primary forms of abstraction have been studied extensively. \emph{State abstraction} compresses the state space by grouping states that are equivalent with respect to some criterion. \citet{li2006towards} provide a unifying framework for state abstraction in MDPs, defining several abstraction types based on what properties they preserve (optimal value, optimal action, model equivalence). \citet{abel2020theory} extends this theory to characterize the tradeoffs between abstraction granularity and value loss, showing that approximate state abstractions can yield near-optimal behavior while substantially reducing the computational burden of learning and planning.

\emph{Action abstraction} (or temporal abstraction) extends the agent's decision-making to include temporally extended actions. The options framework \citep{sutton1999options} formalizes this through the concept of an option $\omega = \langle \mathcal{I}_\omega, \Policy_\omega, \beta_\omega \rangle$, consisting of an initiation set $\mathcal{I}_\omega \subseteq \States$ specifying where the option can be initiated, a policy $\Policy_\omega$ that governs behavior during execution, and a termination condition $\beta_\omega : \States \rightarrow [0,1]$. Options can be incorporated into planning and learning algorithms as temporally extended choices alongside primitive actions, enabling agents to reason and learn at multiple levels of temporal granularity.

A key insight connecting these two forms of abstraction comes from work on grounding symbolic representations in continuous environments. \citet{konidaris2018skills} show that the symbols necessary and sufficient for planning with a given set of skills can be derived from the initiation sets and effect distributions of those skills. This result establishes a principled link between action abstraction (skills or options) and state abstraction (symbolic predicates): the appropriate symbolic vocabulary for planning is determined by the agent's available actions, not imposed externally. \citet{konidaris2016constructing} further develops this into the \emph{skill-symbol loop}, in which acquiring new skills gives rise to new symbolic representations, which in turn enable planning at higher levels of abstraction. Each level of the hierarchy abstracts away details of the levels below, exposing the structure of the task at an appropriate granularity.

The IPT formalism introduced in \cref{sec:hybrid} instantiates these ideas. The detector $\Abstraction : \States \rightarrow \SymStates$ maps subsymbolic states to symbolic descriptions, and the executor mapping $e : \Operators \rightarrow \Executors$ provides action abstraction: each symbolic operator $\Operator$ abstracts over its executor $e(\Operator)$, representing a potentially complex policy as a single named step with well-defined preconditions and effects. What makes these abstractions \emph{structured}---as opposed to, say, a learned latent representation in an end-to-end model---is that they preserve compositional relationships: operators compose into plans, preconditions and effects support causal reasoning, and the abstraction hierarchy creates natural interfaces at which mismatch can be detected and addressed. This compositional structure is what distinguishes the approach taken in this dissertation from approaches that learn monolithic mappings from observations to actions \citep{konidaris2019necessity}.

\subsection{Action Abstractions and Adaptation}
\label{sec:action_abstractions}

The first form of structured abstraction we study concerns the representation of actions. In the IPT formalism (\cref{sec:hybrid}), the agent's behavior is organized as a sequence of operators, each associated with an executor. This decomposition creates a structured relationship between planning and execution that can be exploited when novelty causes failure.

When an executor $e(\Operator)$ fails---because the environment dynamics have changed in a way that the executor's policy can no longer achieve $\Eff_{\Operator}$---the failure can be localized to a specific point in the plan. The agent knows which operator was being executed, what its intended preconditions and effects were, and what actually occurred, assuming the detector can report whether the expected symbolic effects were achieved. This information converts what would otherwise be a global performance degradation into a targeted learning problem: the agent needs to acquire or revise an executor for a specific operator, with initiation and termination conditions supplied by the symbolic specification.

This targeted formulation has two important consequences. First, it restricts the scope of learning: the agent does not need to treat the entire task as a new learning problem, but can focus exploration on the component whose expected effects failed. This leads to substantially improved sample efficiency compared to unstructured learning approaches in the settings studied in \cref{ch:rapid_learn}. Second, it preserves previously acquired knowledge: executors for operators that were not affected by the novelty remain valid and can be reused. This connects to the broader theme in the abstraction literature that well-chosen abstractions support knowledge retention and transfer \citep{abel2018state, konidaris2007building}.

The effectiveness of this approach depends directly on the quality of the action abstraction. If operators are too coarse, failures are difficult to localize. If operators are too fine, the advantage of symbolic reasoning is lost. The appropriate level of abstraction is one where operators correspond to semantically meaningful units of behavior whose preconditions and effects capture the causally relevant aspects of the interaction---precisely the kind of abstraction that \citet{konidaris2018skills} show is necessary and sufficient for sound planning. We develop this approach concretely in \cref{ch:rapid_learn}, presenting a hybrid planning--learning method that leverages action abstractions for efficient adaptation under novelty.

\subsection{Interaction Abstractions and Generalization}
\label{sec:interaction_abstractions}

The second form of structured abstraction concerns the representation of interaction between the agent and objects in the environment. While action abstractions organize behavior into compositional sequences, interaction abstractions determine \emph{how} individual actions are represented---specifically, what features of the agent--object interaction are captured by the policy's state and action spaces.

In robotic manipulation, many robot-learning approaches define policies over robot- or controller-specific representations: joint angles, end-effector poses, or joint torques. These representations can tightly couple the learned behavior to a specific robot embodiment and control interface. The abstraction literature offers a lens for understanding this limitation: a robot-centric representation may fail to abstract away the details of the agent, and therefore may not capture the natural complexity of the task \citep{konidaris2019necessity}. As a result, a policy trained on one robot often cannot be applied directly to another, even when both robots are performing the same physical task on the same object. Similarly, learned behavior may fail to generalize across objects whose task-relevant physical properties differ from the training instance.

An alternative is to represent interaction at the level of the object rather than the robot. Object-centric force-space representations describe what forces are applied to the object, where on the object they are applied, and what motion the object undergoes as a result. This shift constitutes a structured interaction abstraction: it organizes the policy's state and action spaces around the physically meaningful invariants of the task---the forces and motions that determine success---rather than the robot-specific details that vary across platforms.

This form of abstraction is analogous to the portable option representations studied by \citet{konidaris2007building}, where temporally extended skills are defined to transfer across task instances. Object-centric representations more broadly seek to factor manipulation knowledge away from specific robots and scenes \citep{kroemer2021review,james2022autonomous}. In our case, the ``portability'' arises not from learning option policies intended to transfer wholesale, but from choosing a representation space---forces and object geometry---in which the physics of the interaction is invariant across the variations that matter (different objects, different robots, different physical parameters).

The consequence is a more principled route to generalization. A policy learned in force space over one object is better positioned to generalize to other objects that share similar physical semantics, because the representation captures aspects of the interaction that are invariant across objects rather than those specific to a single training instance. Similarly, a policy learned on one robot can be executed on a different robot when the new robot can realize the required forces at the required contact points. We develop this approach in \cref{ch:flex,ch:pho} through force-based learning methods for robotic manipulation that generalize across objects, physical parameters, and robotic embodiments.

% Structured abstractions conceptual figure
\begin{figure}[t]
    \centering
    \begin{tikzpicture}[
        >=Stealth,
        % --- Node styles (matching Ch.4 palette) ---
        srcred/.style={
            rectangle, rounded corners=3pt, draw=red!60!black, fill=red!8,
            minimum width=3.0cm, minimum height=0.7cm,
            font=\small, align=center, thick
        },
        srcblue/.style={
            rectangle, rounded corners=3pt, draw=blue!60!black, fill=blue!8,
            minimum width=3.0cm, minimum height=0.7cm,
            font=\small, align=center, thick
        },
        absbox/.style={
            rectangle, rounded corners=3pt, draw=purple!60!black, fill=purple!8,
            minimum width=3.4cm, minimum height=1.0cm,
            font=\small, align=center, thick
        },
        capbox/.style={
            rectangle, rounded corners=3pt, draw=green!60!black, fill=green!8,
            minimum width=3.4cm, minimum height=1.0cm,
            font=\small, align=center, thick
        },
        % --- Arrow style ---
        flowarrow/.style={->, thick, >=Stealth},
        % --- Labels ---
        collabel/.style={font=\small\bfseries, gray!60!black},
        chaplabel/.style={font=\scriptsize\itshape, gray!45},
        % --- Legend ---
        legendbox/.style={minimum width=0.45cm, minimum height=0.3cm, rounded corners=2pt, thick, inner sep=0pt},
        legendlabel/.style={font=\scriptsize, anchor=west, inner sep=0pt},
    ]

    % ===== Column headers =====
    \node[collabel] at (-4.8, 0.8) {Source};
    \node[collabel] at (0.0, 0.8) {Abstraction};
    \node[collabel] at (4.8, 0.8) {Capability};
    \draw[gray!30, thin] (-6.8, 0.5) -- (6.8, 0.5);

    % =========================================================
    %  ROW 1: Novelty → Action Abstractions → Adaptation
    % =========================================================

    \node[srcred] (n1) at (-4.8, -0.2) {Changed dynamics};
    \node[srcred] (n2) at (-4.8, -1.2) {Execution failures};

    \node[absbox] (aa) at (0.0, -0.7) {\textbf{Action abstractions}\\[-2pt]{\scriptsize Operators $\leftrightarrow$ executors}};

    \node[capbox] (ad) at (4.8, -0.7) {\textbf{Efficient adaptation}\\[-2pt]{\scriptsize Failure localization}\\[-2pt]{\scriptsize Targeted learning}};

    \draw[flowarrow, red!50!black]   (n1.east) -- (aa.west |- n1);
    \draw[flowarrow, red!50!black]   (n2.east) -- (aa.west |- n2);
    \draw[flowarrow, purple!50!black] (aa.east) -- (ad.west);

    \node[chaplabel] at (4.8, -1.6) {Chapter~4};

    % =========================================================
    %  Separator
    % =========================================================
    \draw[gray!25, thin, densely dashed] (-6.8, -2.2) -- (6.8, -2.2);

    % =========================================================
    %  ROW 2: Variation → Interaction Abstractions → Generalization
    % =========================================================

    \node[srcblue] (v1) at (-4.8, -2.9) {Different objects};
    \node[srcblue] (v2) at (-4.8, -3.9) {Different embodiments};

    \node[absbox] (ia) at (0.0, -3.4) {\textbf{Interaction abstractions}\\[-2pt]{\scriptsize Force-space representations}};

    \node[capbox] (ge) at (4.8, -3.4) {\textbf{Generalization}\\[-2pt]{\scriptsize Cross-object transfer}\\[-2pt]{\scriptsize Cross-embodiment transfer}};

    \draw[flowarrow, blue!50!black]   (v1.east) -- (ia.west |- v1);
    \draw[flowarrow, blue!50!black]   (v2.east) -- (ia.west |- v2);
    \draw[flowarrow, purple!50!black] (ia.east) -- (ge.west);

    \node[chaplabel] at (4.8, -4.3) {Chapter~5};

    % =========================================================
    %  Legend (centered, single row)
    % =========================================================
    \draw[gray!30, thin] (-6.8, -4.9) -- (6.8, -4.9);

    \node[legendbox, draw=red!60!black, fill=red!8]      (L1) at (-4.0, -5.3) {};
    \node[legendlabel] at (-3.7, -5.3) {Novelty};

    \node[legendbox, draw=blue!60!black, fill=blue!8]     (L2) at (-1.8, -5.3) {};
    \node[legendlabel] at (-1.5, -5.3) {Variation};

    \node[legendbox, draw=purple!60!black, fill=purple!8]  (L3) at (0.6, -5.3) {};
    \node[legendlabel] at (0.9, -5.3) {Abstraction};

    \node[legendbox, draw=green!60!black, fill=green!8]    (L4) at (3.2, -5.3) {};
    \node[legendlabel] at (3.5, -5.3) {Capability};

    \end{tikzpicture}
    \caption[Structured abstractions mediating agent responses to novelty and variation]{The role of structured abstractions in mediating the agent's response to novelty and variation. \textbf{Top:} When novelty causes changed dynamics or execution failures, action abstractions (the operator--executor decomposition) enable failure localization and targeted learning, supporting efficient adaptation (Chapter~4). \textbf{Bottom:} When variation arises from different objects or embodiments, interaction abstractions (force-space representations) capture task-relevant physical invariants, supporting generalization without retraining (Chapter~5). In both cases, the abstraction structures the agent's response rather than requiring the agent to relearn from scratch.}
    \label{fig:abstractions_buffer}
\end{figure}

\subsection{The Unifying Perspective}
\label{sec:unifying_perspective}

Both forms of abstraction share a common principle, illustrated in \cref{fig:abstractions_buffer}: they organize the agent's representation around aspects of the task and interaction that are causally relevant, rather than aspects that are contingent on a particular environment instance or robot embodiment.

Action abstractions organize behavior into operators with explicit preconditions and effects, so that when novelty disrupts one operator, the disruption is contained and the agent can focus learning on the affected component. Interaction abstractions organize execution around physical invariants---forces, contact points, object-relative motion, and object dynamics---so that when the specific robot or object changes, the policy's representation can still capture the structure of the task.

In both cases, the abstraction serves as a buffer between the agent and the variability of the world. It does not eliminate novelty or variation---that would require closed-world assumptions---but it structures the agent's response to change. In \cref{ch:rapid_learn}, this structure supports efficient recovery after an execution failure. In \cref{ch:flex,ch:pho}, it supports transfer across anticipated variation before failure occurs. The empirical validation of this perspective is the subject of the remaining chapters.
